% Accepted after independent root and intervals audits and integrated in sections/C.tex, 18 November 2026.
\begin{claim}{C.27}{Proved under stated hypotheses}{A flat interval can be approximated by split inclusions}\label{C27}
Assume the corollary setting of Claim~\ref{C6}, including the canonical normalization and full interval hypothesis. Choose a presentation $\mathcal M\simeq\operatorname{Fun}_k(\mathcal A^{\mathrm{op}},\Mod k)$ by a small $k$-linear category $\mathcal A$ with projective Hom modules. We do not require this indexing category to contain zero objects or finite direct sums; its additive envelope supplies the additive generator category of Claim~\ref{C6}. Suppose the monoidal structure is pointwise on this chosen presentation: for $a\in\mathcal A$ we have specified natural identifications
\[
 \mathbf1(a)=k,\qquad (X\otimes Y)(a)=X(a)\otimes_kY(a)
\]
compatible with the unit constraints. Suppose also that the interval object has a filtered presentation
\[
 J\cong\mathop{\mathrm{colim}}_{i\in I}P_i
\]
by finitely generated projective objects of $\mathcal M$, and that $\alpha:\mathbf1\to J$ is represented by a map $\alpha_{i_0}:\mathbf1\to P_{i_0}$ at one stage. Then every projective resolving adjunction $F\dashv G$ of Claim~\ref{C6} has a raw right adjoint inducing the equivalence of the manuscript's corollary.

The factorization hypothesis holds, in particular, if $\mathbf1$ is finitely presented in $\mathcal M$ and $J$ is flat as a module over the small linear category, with flatness expressed by such a filtered presentation. Neither finite presentation of the unit nor flatness of the interval is asserted to follow from the manuscript.
\end{claim}
\begin{proof}
We may replace $I$ by the filtered undercategory $i_0/I$. Its forgetful functor to $I$ is final, so the colimit remains $J$, and every stage carries a specified compatible map $\alpha_i:\mathbf1\to P_i$ whose composite to $J$ is $\alpha$. For every $X$, the map $\alpha_i\otimes X$ is monic: its composite with $P_i\otimes X\to J\otimes X$ is the monomorphism $\alpha\otimes X$.

We first check that $\alpha_i(a):k\to P_i(a)$ splits as a $k$-module map. Apply the preceding monicity to the scalar copower $\kappa(V)=V\otimes_k\mathbf1$, for every $k$-module $V$. The pointwise tensor formula identifies the resulting map at $a$ with $\alpha_i(a)\otimes_kV$. Thus $\alpha_i(a)$ is pure. The module $P_i(a)$ is projective over $k$: a finitely generated projective in the functor category is a retract of a finite sum of representables, whose evaluations are projective by the corollary's Hom-projectivity assumption.

For completeness, a pure map $u:k\to P$ into a projective $k$-module splits even when $P$ is not finitely generated. Embed $P$ as a direct summand of a free module $k^{(S)}$. The image of $u(1)$ has finite support, with coordinates $c_1,\ldots,c_m$. The composite $k\to k^{(S)}$ remains pure. If $L=(c_1,\ldots,c_m)$, tensoring with $k/L$ sends this composite to the zero map, so purity forces $k/L=0$. Choose $r_1,\ldots,r_m$ with $\sum_jr_jc_j=1$. The corresponding finite-coordinate functional $k^{(S)}\to k$, restricted to $P$, retracts $u$.

Fix such retractions of $\alpha_i(a)$ for all $a$. They need not be natural in $a$. The pointwise tensor formula nevertheless shows that every evaluation of $\alpha_i\otimes X$ is a split monomorphism of $k$-modules. Since every component $F_n$ is a retract of a sum of evaluation functors, each degree of $F(\alpha_i\otimes X)$ is split monic as well.

Let $K$ be an acyclic nonnegative complex of projective $k$-modules, and put $Q=GK$. By Claim~\ref{C1}, $K$ is a degreewise finite product of disks $D^m(Q_m)$ with projective coefficients $Q_m$. Such a complex is injective relative to degreewise split monomorphisms of complexes. Indeed a chain map into $D^m(Q_m)$ is specified by a module map in degree $m-1$; a splitting in that degree extends it, and the product of the extensions is a chain map into $K$.

We can therefore extend the counit $\epsilon_K:FQ\to K$ across every $F(\alpha_i\otimes Q)$ to a chain map
\[
 e_i:F(P_i\otimes Q)\longrightarrow K.
\]
Its adjunct $r_i:P_i\otimes Q\to GK=Q$ satisfies $r_i(\alpha_i\otimes Q)=\id_Q$, because the adjunct of the counit is $\id_{GK}$. Thus $Q$ is a retract of $P_i\otimes Q$ for each stage. We do not require the retractions $r_i$ to be compatible, and we do not assert that the individual objects $P_i\otimes Q$ are null.

The pointwise tensor and all components of $F$ preserve filtered colimits. Consequently
\[
 \mathop{\mathrm{colim}}_i F(P_i\otimes Q)
 \cong F(J\otimes Q)\simeq0.
\]
Filtered colimits of modules are exact, so the same equality holds after taking homology in any fixed degree. Let $z\in H_n(FQ)$, and map it to a class $z_i\in H_nF(P_i\otimes Q)$ at one stage. Its image in the displayed colimit is zero. The elementary description of a filtered colimit of modules supplies a later stage $i\to j$ at which the image of $z_i$ is zero. Compatibility of the maps $\alpha_i$ says that this image is $H_nF(\alpha_j\otimes Q)(z)$. Apply $H_nF(r_j)$ to conclude $z=0$.

We have proved $FGK\simeq0$ for every acyclic nonnegative projective complex $K$, exactly the remaining test of Claim~\ref{C10}. Canonical normalization proves the degree-zero projective-coefficient counit, as in Claim~\ref{C7}. Claims~\ref{C10} and~\ref{C3} now prove the asserted raw equivalence for the originally chosen resolving adjunction.
\end{proof}
